
\section{Novikov isocrytals}

In this section we define the notion of a Novikov isocrystal and prove some
basic facts about it. Throughout the discussion, we fix a real number $\lambda >
1$.

\subsection{Novikov isocrystals}

\begin{definition}
  \label{def:frobenius}
  \leanok
  \uses{def:novikov_series}
  For an abelian group $A$, we define the Frobenius endomorphism $F \colon
  A((t^{\mathbb{R}})) \to A((t^{\mathbb{R}}))$ by $A$-linearly extending $t^d
  \mapsto t^{\lambda d}$.
\end{definition}

\begin{theorem}
  \label{thm:frobenius_properties}
  \lean{Novikov.frobeniusRingHom, Novikov.frobenius_continuous}
  \leanok
  \uses{def:frobenius, def:novikov_ring, thm:novikov_topology}
  When $A$ is a ring, the Frobenius is a continuous ring homomorphism.
\end{theorem}

\begin{proof}
  \leanok
  Because both sides of the convolution variables are equally scaled, it
  respects multiplication. It also sends the standard neighborhood bases to
  neighborhood bases.
\end{proof}

\begin{definition}[Definition~3.1]
  \label{def:novikov_isocrystal}
  \lean{Novikov.NovikovIsocrystal}
  \leanok
  \uses{def:frobenius, thm:novikov_ring, thm:frobenius_properties}
  Let $A$ be a commutative ring. A \emph{Novikov isocrystal} over $A$ consists of
  \begin{itemize}
    \item a finite projective $A((t^{\mathbb{R}}))$-module $M$, and
    \item a $F$-semilinear isomorphism $F_M \colon M \to M$.
  \end{itemize}
\end{definition}

\begin{lemma}
  \label{lem:novikov_isocrystal_category}
  \leanok
  \uses{def:novikov_isocrystal, def:novikov_isocrystal_hom}
  Novikov isocrystals over $A$ form a preadditive category with internal Hom.
\end{lemma}

\begin{proof}
  \leanok
  Given two isocrystals $M$ and $N$, the module
  $\operatorname{Hom}_{A((t^{\mathbb{R}}))}(M, N)$ is finite projective and
  inherits a Frobenius action from those on $M$ and $N$.
\end{proof}

\begin{theorem}[Proposition~3.3]
  \label{thm:vect_to_isoc_full_faith}
  \lean{Novikov.NovikovIsocrystal.vectToNovIsoc_fully_faithful}
  \leanok
  \uses{def:novikov_isocrystal}
  The functor $\mathsf{Vect}(A) \to \mathsf{NovIsoc}(A)$ sends a finite
  projective $A$-module $M$ to its constant isocrystal $M((t^{\mathbb{R}}))
  = M \otimes_A A((t^{\mathbb{R}}))$ is fully faithful.
\end{theorem}

\begin{proof}
  \leanok
  A morphism $g$ of constant isocrystals corresponds to an
  $A((t^{\mathbb{R}}))$-linear map $M \otimes_A A((t^{\mathbb{R}})) \to
  N \otimes_A A((t^{\mathbb{R}}))$ commuting with $\mathrm{id} \otimes F$. Via
  the internal Hom, this yields an element of $A((t^{\mathbb{R}})) \otimes_A
  \operatorname{Hom}_A(M, N)$ fixed by $F$. Such elements are exactly of the
  form $1 \otimes f$ for $f \in \operatorname{Hom}_A(M, N)$.
\end{proof}

\subsection{Some lemmas}

\begin{theorem}[Proposition~3.5]
  \label{thm:frobenius_eigenvector}
  \lean{Novikov.exists_frobenius_eigenvector}
  \leanok
  \uses{def:novikov_isocrystal, def:frobenius, thm:novikov_field, thm:novikov_topology,
    def:novikov_filtration, def:novikov_series}
  Let $K$ be an algebraically closed field and $M$ a nonzero Novikov
  isocrystal over $K$. Then there exists a nonzero $m \in M$ and
  $c \in K^\times$ such that $F_M(m) = c \cdot m$.
\end{theorem}

\begin{proof}
  \leanok
  Pick a nonzero $m_0 \in M$ and choose the minimal $n \ge 1$ such that
  $F_M^0 m_0, \dots, F_M^{n-1} m_0$ are linearly independent while
  $F_M^n m_0 = \sum_{i} a_i \cdot F_M^i m_0$ for some series $a_i$.
  Scale $m_0$ by a monomial $t^r$ so that the new coefficients $a'_i$
  lie in $K[\![t]\!]$ and at least one has nonzero constant term
  $\bar{a}_i$.  By algebraic closedness, find $c \in K^\times$ with
  $c^n = \sum_i \bar{a}_i c^i$.  Define the contraction
  $
    g(x) = \sum_{i=0}^{n-1} c^{-(n-i)} F^{n-i}(x) F^{n-1-i}(a'_i).
  $
  If $x - y$ lies in the $D$-filtration, then $g(x) - g(y)$ lies in the
  $(\lambda \cdot D)$-filtration, and $g(1)$ is in the $0$-filtration
  with $(g(1))_0 = 1$.  Iterating $g$ from $1$ produces a Cauchy sequence
  converging to a fixed point $b$ with $b_0 = 1$.  From $b$ define
  $b_0 = c^{-1}F(b)a'_0$ and
  $b_{k+1} = c^{-1}(F(b_k) + F(b)a'_{k+1})$; then $b_{n-1} = b$ and
  $
    m' = \sum_{i=0}^{n-1} b_i \cdot F_M^i m_0
  $
  is a nonzero eigenvector satisfying $F_M(m') = c \cdot m'$.
\end{proof}

